% MT005 / explanation — canonical editable content.
先把三个概率分清。\(P(y=c)\)是\textbf{先验}，表示还没看到这条输入时，类别\(c\)有多常见；\(p(x\mid y=c)\)是\textbf{类条件分布}，表示已知是\(c\)类时，什么输入常见；\(P(y=c\mid x)\)是\textbf{后验}，表示看到输入\(x\)以后，它属于\(c\)类的概率。竖线“\(\mid\)”可读成“在已知……的条件下”。


NB先学习前两项，再用Bayes公式得到后一项。比如先统计垃圾邮件的比例，以及垃圾邮件和正常邮件分别常用哪些词，再判断一封新邮件。每个选项都应放回这条流程理解：



\begin{itemize}

\item \textbf{A按原解选。} 训练最终使我们能够估计类别后验，但生成式模型直接学习的是先验与类条件分布。这里“等价”指最终作用，不能理解成直接用后验作训练模型；\StudyRef{Answers}{mt032}{MT032}的措辞差异也在这里。

\item \textbf{B按原解选。} 决策边界附近两类势均力敌，模型应表现出不确定性。把标签写成0或1，若\(P(Y=1\mid x)=p\)，则标签方差为\(p(1-p)\)：\(p=0.5\)时为0.25，\(p=0.9\)时为0.09。这里是标签的不确定性，不是反复训练所得参数的波动。

\item \textbf{C选。} 非高斯的类条件分布可以产生非线性边界。这只是“可以”，没有说只有非高斯才可以；各类方差不同的Gaussian NB也常有弯曲边界。

\item \textbf{D选。} 最大似然估计（maximum likelihood estimation，MLE）选择使已观察数据最可信的参数。NB拟合各类的数据分布，LR拟合给定输入后的标签概率；两者的学习对象不同，但都能使用MLE。

\item \textbf{E不选。} NB也能正则化。例如某词在一类中碰巧没出现，直接估计概率为0会使整条概率乘积归零。加平滑给各事件少量额外计数，可以缓和小样本导致的极端估计。

\end{itemize}


选读：二值标签满足\(Y^2=Y\)，所以\(E[Y^2]=E[Y]=p\)。按方差定义，\(\operatorname{Var}(Y)=E[Y^2]-E[Y]^2=p-p^2\)，这就是B中公式的来源。


\textbf{出处：} 2020B期中Quiz Q5，原卷第3页（5分）


\textbf{答案依据：} 2020B Quiz随卷解答 / 2020B随卷答案副本 Q5，随卷答案ABCD；解释A/B语境。
